\chapter{Problem Statement}

\hspace{\parindent} The public transportation system has long been a commodity that comes in the form of bus transport, train networks, and the like, with predetermined routes and fares \cite{EdsaCarousel2025Overview}. In Metro Manila, traffic congestion has long been a problem stemming from insufficient road infrastructure and poor traffic mitigation policies, resulting in an inefficient public transport system \cite{NEEDCITE_congestion}. One of the government's actions for an alternative transportation system that runs within the metropolitan area is the implementation of the EDSA Carousel in 2020. It is the only busway system in Metro Manila implemented by the government to have a designated lane, determined by concrete barriers \cite{EdsaCarousel2025Overview}. Over the years, the number of commuters using the EDSA Carousel has steadily risen, averaging approximately $1.83\times10^5$ daily passengers, with single-day peaks exceeding $3.2\times10^5$ during high-demand periods, reaching a new high of nearly 67 million in 2025, up from 63.02 million in 2024 \cite{DOTr2025Ridership}. However, despite its growing demand, the system still faces several problems that can affect the quality of bus routing. To start, the LTFRB relied on a tedious manual monitoring process by deploying personnel at each stop \cite{NEEDCITE_manualmonitoring}. There are also certain portions where it operates with mixed traffic, causing delays among buses. It is also clear that bus bunching has not been resolved; most of the time, their designated dispatch schedule is not followed \cite{NEEDCITE_busbunching_edsa}. Moreover, the operation of the bus system is also largely influenced by external factors, such as weather and road maintenance. Sections of the EDSA Bus Carousel routes typically close during tropical cyclones to prevent potential damage from continuous rain and strong winds \cite{DOTr2020Suspension}. Given these operational issues, the bus dispatching process of the EDSA Bus Carousel can be modeled as a Vehicle Scheduling Problem (VSP), in which buses must be dispatched and coordinated efficiently despite these issues. By improving the scheduling, operational issues, such as bus bunching, irregular dispatch intervals, and passenger waiting times, will be reduced, thereby improving the overall bus dispatch efficiency.

Busway systems operate in a highly dynamic, unpredictable environment. In the past, Vehicle Scheduling Problems (VSP) relied on traditional scheduling techniques that focused on finding the optimal path by determining the shortest distance. Static scheduling is the pre-computed timetabling that fixes departure times before operation begins, with no mechanism to revise them in response to real-time conditions \cite{Ju2023Joint,Cao2022Train,Nie2025CMRM}. However, because transit environments are inherently stochastic, these static schedules perform poorly during daily operations, which causes severe operational issues such as bus bunching, node failures, and an increase in passenger wait times \cite{Daganzo2009,FEMbusbunching}.

To address the rigidity of static schedules, recent scheduling studies have shifted toward dynamic, AI-driven approaches, with Multi-Agent Reinforcement Learning (MARL) as the leading solution. Past studies have successfully utilized MARL to overcome the scalability limits of earlier systems by distributing the scheduling decisions across autonomous agents, improving the overall system's adaptability to real-time traffic \cite{Rodriguez2023Cooperative,Wang2023MultiObj,Wangsun}. However, a significant gap remains in how these state-of-the-art models are evaluated. Existing literature primarily validates the MARL scheduling models in a controlled, simulated environment. Previous evaluations rarely account for the heavy-tailed, asymmetric disruptions characteristic of actual transit operations, leaving their performance under non-ideal, real-world traffic anomalies largely untested. To the best of our knowledge, no existing MARL bus-scheduling model has been sufficiently tested against severe, simultaneous disruptions such as sudden passenger surges and weather-induced delays. To bridge this gap, this study develops and evaluates a MARL-based bus scheduling model within a stochastic traffic simulation subjected to non-ideal conditions.

This study develops and evaluates a MARL-based bus scheduling model on a calibrated stochastic simulation of the EDSA Carousel northbound corridor, with non-ideal conditions operationalized as the simultaneous and independent injection of passenger demand surges, severe-weather travel-time skew, and Poisson-distributed bus breakdowns. The objective is to quantify how MARL performance degrades from its ideal-condition baseline and how quickly it recovers, producing the first integrated robustness characterization for MARL bus scheduling under this specific combination of disturbances.

\section{Rationale}
Urban bus corridors such as the EDSA Carousel handle hundreds of thousands of daily passengers and operate under conditions that routinely violate the assumptions of existing MARL simulations: passenger volumes surge during peak periods and weather events, travel times skew sharply during typhoons and heavy rainfall, and individual buses experience mechanical failures during service. The scale of the corridor makes these disturbances operationally consequential: the EDSA Carousel averaged roughly $1.83\times10^5$ passengers per day in 2025 with single-day peaks exceeding $3.2\times10^5$ \cite{DOTr2025Ridership}, and its operations are suspended outright under elevated tropical-cyclone signals \cite{DOTr2020Suspension}, so weather effects manifest both as heavy-tailed travel-time delays on operating days and as complete service withdrawal on the most severe days. If MARL is to be considered a credible candidate for deployment in such corridors, its performance must be quantified under exactly these disturbances. Reporting wait-time and headway-regularity results only under stationary, idealized conditions overstates expected operational benefit and provides transit operators no basis for risk assessment.

For clarity, a \textit{stochastic environment} in this study models uncertain quantities as random variables drawn from calibrated probability distributions rather than as fixed values. Each Monte Carlo run produces a different realization; the simulator captures the shape of uncertainty (mean, variance, and tail behavior) rather than a single predetermined trajectory. Modeling something as a random variable is not the same as making it deterministic. A coin flip remains stochastic even when modeled as $\text{Bernoulli}(0.5)$, because the outcome of any single flip is not predictable from the distribution.

Addressing the robustness gap is timely because (a) recent MARL bus-control work has explicitly identified robustness under perturbations as the next research priority \cite{Wangsun}, (b) microsimulation tools such as SUMO now support empirical calibration to real corridor data via standard GEH and RMSE metrics \cite{FHWA2019TAT,Patil2025Conformal}, and (c) calibrated stochastic models of weather-induced and demand-driven disturbances are now established in the transportation forecasting literature \cite{Patil2025Conformal}. Together, these make an integrated robustness evaluation feasible that was not previously possible.

\section{Research Gap}
Existing Multi-Agent Reinforcement Learning frameworks for bus scheduling demonstrate measurable improvements in passenger waiting time and headway regularity, but these gains are validated almost exclusively in simulation environments that assume stationary passenger arrivals, symmetric Gaussian travel times, and a failure-free fleet. The performance of MARL bus-scheduling models has not been characterized when these idealizing assumptions are simultaneously violated by (i) stochastic passenger demand surges, (ii) heavy-tailed weather-induced travel-time delays, and (iii) discrete bus breakdowns. Without this characterization, it cannot be determined whether reported MARL gains persist, degrade gracefully, or collapse under realistic operating disturbances, which in turn blocks the transition of MARL bus scheduling from simulation to real urban transit deployment.



\section{Research Objectives}
\subsection{Main Objective}
To investigate the performance of a developed Multi-Agent Reinforcement Learning (MARL) scheduling model for dynamic bus scheduling under non-ideal traffic conditions on a calibrated simulation of the EDSA Carousel northbound corridor.

\subsection{Specific Objectives}
\begin{enumerate}
   \item \textbf{Environment Construction:} To construct a stochastic traffic simulation environment representing a defined traffic area.
        \\
        \\ \textbf{Expected Output 1.1:} A simulation model of a defined operational scope, validated against empirical traffic data to establish an environment that accurately approximates real-world traffic.
        \\
        \\ \textbf{Expected Output 1.2:} An environment simulator capable of introducing environmental variability through adjustable parameters to simulate non-ideal operating conditions.
        \\
   \item \textbf{Algorithm Implementation:} To develop and implement a Multi-Agent Reinforcement Learning (MARL) approach for dynamic bus scheduling.
        \\
        \\ \textbf{Expected Output 2.1:} A formulated    agent architecture defining individual bus units as agents with discrete state and action space with a continuous reward function that punishes irregularities and passenger service delays.
        \\
   \item \textbf{Performance Evaluation Under Ideal and Non-Ideal Conditions:} To evaluate the scheduling performance of the MARL approach under both ideal operating conditions and non-ideal conditions.
        \\
        \\ \textbf{Expected Output 3.1:} A quantitative assessment of passenger waiting time and travel times under ideal conditions, comparing the MARL approach against baseline scheduling methods.
        \\
        \\ \textbf{Expected Output 3.2:} A statistical analysis of arrival consistency under ideal conditions, using headway variability to determine the system's effectiveness in maintaining uniform gaps between vehicles.
        \\
        \\ \textbf{Expected Output 3.3:} A quantitative assessment of passenger waiting time and travel times under non-ideal conditions, measured as degradation relative to the ideal-conditions performance.
        \\
        \\ \textbf{Expected Output 3.4:} A statistical analysis of the impact of non-ideal conditions, using headway variability to determine the system's effectiveness in maintaining uniform gaps between vehicles under disturbance.
        \\
 \end{enumerate}


\section{Significance of the Study}
This study contributes both practical and scientific significance.

\textbf{Practical significance.} The EDSA Carousel carries on the order of $1.8\times10^5$ passengers daily, with single-day peaks exceeding $3.2\times10^5$ \cite{DOTr2025Ridership}. Even a 30-second reduction in mean passenger waiting time at the average-day scale translates to roughly $1.5\times10^3$ passenger-hours saved per day. This study also quantifies the \textit{recovery time} after a disruption, defined as the simulation time required for headways to re-stabilize after a breakdown, surge, or weather event. This metric is not reported in prior MARL bus-scheduling work but is directly relevant to service reliability and operator decision-making during disruptions.

\textbf{Scientific significance.} The study establishes a reproducible stress-testing protocol for bus-scheduling reinforcement learning: a calibrated corridor model paired with three configurable stochastic anomaly generators (demand, weather, breakdown). Future MARL work can be benchmarked on the same protocol, allowing performance claims to be compared on a common basis. The study also reports the performance-degradation curve of MARL as a function of anomaly intensity, characterizing robustness rather than only average-case performance. Without this characterization, MARL bus-scheduling models risk being deployed on the basis of ideal-simulator benchmarks, with their real-world failure modes unknown until deployment.

\section{Scope and Limitations}

\textbf{Operational definition of non-ideal conditions.} For this study, non-ideal conditions are operationally defined as operating regimes in which one or more of the following four disturbance classes are active: (i) passenger demand surges modeled as a multiplicative scaling factor drawn from $\mathcal{N}(1, \sigma_d^2)$ clipped to $[1, 3]$; (ii) traffic-speed perturbations modeled as a multiplicative scaling factor drawn from $\mathcal{N}(1, \sigma_s^2)$ clipped to $[0.8, 1.2]$; (iii) severe-weather-induced inter-stop travel times sampled from $\text{LogNormal}(\mu_{ln}, \sigma_{ln})$ where the coefficient of variation is set by the swept disturbance-intensity parameter $\eta$ (Chapter~3, Section~3.3.3); and (iv) discrete bus breakdowns sampled from a Poisson process with a configurable malfunction rate. All subsequent references to ``non-ideal conditions'' in this manuscript refer to this definition.

The asymmetric demand clip at $[1, 3]$ reflects this study's focus on demand \textit{surges} (the operationally adverse regime for headway stability) rather than symmetric demand variation. The traffic-speed clip at $[0.8, 1.2]$ corresponds to the $\pm 20\%$ daily-congestion friction range used in the perturbation framework adopted from Wang and Sun (Chapter~3, Section~3.3.1) \cite{Wangsun}. Lognormal weather travel-time parameterization follows the precedent established in Chapter~2, which validated lognormal against INRIX freeway data via Kolmogorov-Smirnov test \cite{Patil2025Conformal}. Extreme precipitation events may exhibit heavier tails (Pareto-type) than lognormal can capture; this is treated as a known limitation.

\textbf{Scope.} The study develops and evaluates a MARL bus scheduling framework on a SUMO microsimulation of the EDSA Carousel northbound route. Performance is evaluated across two regimes (ideal and non-ideal) using passenger waiting time, total travel time, headway coefficient of variation, and post-disturbance recovery time as primary metrics, benchmarked against No Control, Forward Headway, and Even Headway holding strategies over $N \geq 30$ Monte Carlo runs per regime.

\textbf{Delimitations.} (a) Due to computational constraints, the simulation is restricted to a defined operational sub-segment of the EDSA Carousel northbound corridor rather than the entire metropolitan road network. The restriction is justified by the need to preserve 1:1 empirical traffic volumes for GEH calibration without resorting to flow scaling. (b) MARL control actions are applied only at designated control stops, selected following criteria established in the literature (see Chapter~2). (c) The study is entirely simulation-based; no on-road deployment is performed. (d) Weather and breakdown effects are modeled as statistical distributions of their impact on travel time and fleet availability, not as physical or mechanical simulations of rain dynamics or vehicle subsystems. (e) Only the four disturbance modes listed above are considered; other disturbances (road accidents, signal failures, driver behavior variability) are out of scope. (f) Sim-to-real transfer is out of scope; the experimental claims are bounded to the calibrated SUMO microsimulation.